\documentclass[conference]{IEEEtran}
\IEEEoverridecommandlockouts
% The preceding line is only needed to identify funding in the first footnote. If that is unneeded, please comment it out.
\usepackage{cite}
\usepackage{amsmath,amssymb,amsfonts}
\usepackage{algorithmic}
\usepackage{graphicx}
\usepackage{textcomp}
\usepackage{xcolor}
\def\BibTeX{{\rm B\kern-.05em{\sc i\kern-.025em b}\kern-.08em
    T\kern-.1667em\lower.7ex\hbox{E}\kern-.125emX}}
\begin{document}

\title{
\thanks{Identify applicable funding agency here. If none, delete this.}
}

\author{\IEEEauthorblockN{1\textsuperscript{st} Given Name Surname}
\IEEEauthorblockA{\textit{dept. name of organization (of Aff.)} \\
\textit{name of organization (of Aff.)}\\
City, Country \\
email address or ORCID}
\and
\IEEEauthorblockN{2\textsuperscript{nd} Given Name Surname}
\IEEEauthorblockA{\textit{dept. name of organization (of Aff.)} \\
\textit{name of organization (of Aff.)}\\
City, Country \\
email address or ORCID}}

\maketitle

\begin{abstract}
This survey paper provides a comprehensive overview of Multimodal Large Language Models (MLLMs), a transformative technology that integrates and processes multiple data modalities, including text, images, audio, and video. The paper explores the background and foundations of MLLMs, their architectures and techniques, vision-language pretraining, and masked language modeling. It also delves into the applications of MLLMs in various domains, highlighting their potential to overcome challenges such as data heterogeneity and domain knowledge sophistication.

The paper identifies key challenges and limitations of current MLLM systems, including the lack of systematic approaches to development, evaluation, and deployment, as well as issues related to hallucination, toxicity detection, and model compression. It also examines the importance of evaluating MLLMs using comprehensive benchmarks, metrics, and adversarial testing techniques.

The survey highlights recent advances in MLLM research, including technical advances, emerging trends, and real-world applications. It emphasizes the need for responsible innovation, integrating ethical standards into the development and deployment of MLLMs to ensure fairness, transparency, and accountability.

This paper contributes to the field by providing a thorough analysis of the current state of MLLM research, identifying key themes, findings, and limitations. It highlights the scope and potential of MLLMs to transform various domains and emphasizes the importance of continued research and development to address the challenges and limitations associated with these models. The survey concludes by outlining future directions and open challenges in MLLM research, emphasizing the need for more effective evaluation methods, better understanding of cognitive abilities, and responsible innovation. Overall, this paper provides a comprehensive overview of MLLMs, their applications, and their potential to achieve artificial general intelligence.
\end{abstract}

\begin{IEEEkeywords}
large language models, multimodal learning, natural language processing, machine learning, artificial intelligence
\end{IEEEkeywords}


\section{Introduction to Multimodal Large Language Models (MLLMs)}
Introduction to Multimodal Large Language Models (MLLMs)

Multimodal large language models (MLLMs) have emerged as a transformative technology, enabling complex AI systems to integrate and process multiple data modalities, including text, images, audio, and video \cite{b1}. The development of MLLMs has been driven by significant advancements in architectures, training methods, and applications, as highlighted in recent surveys and reviews \cite{b2, b3}. This introduction provides an overview of the key findings and contributions of MLLMs, common themes and patterns, contrasting viewpoints or approaches, technical details and methodologies, and limitations and challenges associated with these models.

The capabilities of MLLMs in cross-modal understanding and generation have been demonstrated in various studies \cite{b4, b5}, showcasing their potential to revolutionize applications such as visual storytelling, accessibility, and code generation. Researchers have made notable progress in exploring the technical features, strengths, and limitations of MLLMs, including their architecture designs \cite{b6}, training strategies \cite{b7}, and evaluation metrics \cite{b8}. The importance of large-scale pretraining in developing effective MLLMs has been consistently emphasized \cite{b9, b10}, highlighting the need for robust evaluation benchmarks to assess their performance across different tasks and modalities.

Several common themes and patterns emerge from the literature on MLLMs, including the emphasis on multimodal integration \cite{b11}, large-scale pretraining \cite{b12}, and the need for robust evaluation metrics \cite{b13}. While there are no drastically contrasting viewpoints, some researchers approach the topic from different angles, such as a data-centric perspective \cite{b14} or a technical architecture-focused approach \cite{b15}. The potential applications of MLLMs have also been explored from an application-oriented perspective \cite{b16}, highlighting their potential impact on various fields.

The technical details and methodologies underlying MLLMs are complex and multifaceted, involving architecture designs such as transformer-based models and multimodal fusion techniques \cite{b17}. Training strategies, including large-scale pretraining, fine-tuning, and few-shot learning, have been explored in various studies \cite{b18, b19}. Evaluation benchmarks and metrics have also been introduced to assess the performance of MLLMs across different tasks and modalities \cite{b20}.

Despite the significant progress made in MLLMs, several limitations and challenges remain, including scalability \cite{b21}, robustness \cite{b22}, cross-modal learning \cite{b23}, and ethical considerations \cite{b24}. Developing and deploying MLLMs at scale remains a significant challenge, and researchers must address the potential vulnerabilities of these models to adversarial attacks, data quality issues, and other robustness concerns. Effective cross-modal learning and integration of multiple data modalities remain open research questions, and the ethical implications of developing and deploying MLLMs must be carefully considered.

In conclusion, this introduction provides a comprehensive overview of the current state of MLLMs, highlighting key themes and developments, implications, and connections. By analyzing the literature on MLLMs, we can gain a deeper understanding of their potential applications, limitations, and challenges, providing a foundation for further research and development in this exciting and rapidly evolving field \cite{b25}.

\section{Background and Foundations of MLLMs}
The Multimodal Large Language Models (MLLMs) have been at the forefront of artificial intelligence research, offering significant advancements in understanding and generating multimodal content. To understand the background and foundations of MLLMs, it is crucial to analyze recent papers that delve into their definitions, properties, benchmarking practices, and limitations. A critical examination of these aspects is provided by papers such as \cite{b1}, which contributes by providing a definition of Large Language Models (LLMs) and critically examines common claims regarding their properties, including emergent properties. Furthermore, surveys like \cite{b2} and \cite{b3} offer comprehensive reviews of benchmarks used for evaluating MLLMs across various domains, including understanding, reasoning, generation, and application, highlighting the diversity in task designs, evaluation metrics, and dataset constructions.

The need for clear definitions and standards is a recurring theme in MLLM research, with papers emphasizing the importance of well-defined models and evaluation methods to advance the field cohesively \cite{b1}. The surveys on benchmarks stress that evaluation should be considered a critical discipline supporting the development of MLLMs, indicating a pattern where thorough assessment is key to progress \cite{b2}, \cite{b3}. Moreover, the variety of domains and tasks covered by MLLM research suggests a broad applicability and versatility of these models. While some papers focus on optimizing LLM training and inference through model and system co-design to make LLMs more efficient and accessible \cite{b4}, others primarily concentrate on the evaluation and application aspects without delving into the optimization of training and inference processes \cite{b2}, \cite{b3}.

A critical examination of assumptions and claims about LLMs, as seen in \cite{b1}, contrasts with more overview-oriented surveys that aim to catalog and summarize the state of MLLM research without questioning its foundations \cite{b2}, \cite{b3}. From a technical standpoint, optimizing LLMs involves considering both the architectural design of the models themselves and the system-level optimizations for training and inference efficiency \cite{b4}. The technical approach also involves systematic reviews of existing benchmarks, categorizing them based on domains and evaluating their task designs and metrics \cite{b2}, \cite{b3}.

Despite advancements, significant challenges remain, including the lack of standardization in defining LLMs and in benchmarking practices, which hinders comparative analysis and cohesive progress \cite{b1}, \cite{b2}. Making LLMs efficient, affordable, and accessible remains a challenge, necessitating continued research into model and system design optimizations \cite{b4}. Current evaluation methods for MLLMs are recognized as having limitations, suggesting a need for innovative approaches to assessment that can fully capture the capabilities and potential of these models \cite{b2}, \cite{b3}. This analysis provides a comprehensive overview of the background and foundations of Multimodal Large Language Models, highlighting key contributions, common themes, contrasting viewpoints, technical details, and the challenges faced by this evolving field. The implications of these developments are far-reaching, with potential connections to various applications and domains, underscoring the need for continued research and innovation in MLLM development and evaluation \cite{b1}, \cite{b2}, \cite{b3}, \cite{b4}.

\section{Architectures and Techniques for MLLM Development}
The development of Multimodal Large Language Models (MLLMs) has been a subject of significant interest in recent years, with various architectures and techniques being proposed to improve their performance and applicability. A key finding from the analyzed papers is the introduction of an adaptive self-improvement agentic system that enhances the performance of LLMs in generating ML libraries using architecture-specific programming languages (ASPLs) \cite{b1}. This approach highlights the potential for autonomous improvement of MLLMs, which can lead to more efficient and effective development processes. Furthermore, the application of Technology Readiness Levels (TRL) frameworks has been proposed as a means to ensure robust, reliable, and responsible MLLM development \cite{b2}, providing a principled process for evaluating and improving the maturity of MLLM technologies.

The importance of software architecture in MLLM development is also emphasized, with the identification of four key areas that require attention from both ML and software practitioners to define standard practices for architecting ML-based software systems \cite{b3}. This highlights the need for interdisciplinary approaches that combine expertise from ML, software engineering, and other fields. Additionally, a comprehensive survey of multimodal benchmarks used to evaluate MLLMs has been conducted, highlighting the need for a thorough analysis of benchmarking practices \cite{b4}. A detailed guide to multimodal large language models in vision-language tasks has also been provided, covering essential topics such as training methods, architectural components, and practical applications \cite{b5}.

A common theme across the papers is the need for standardization in MLLM development, with a lack of standardized practices and benchmarks being identified as a significant challenge \cite{b1, b3, b4}. The importance of software architecture and the challenges in developing MLLMs, such as scalability, robustness, and cross-modal learning, are also acknowledged by several authors \cite{b1, b4, b5}. The use of adaptive self-improvement agentic systems, TRL frameworks, and software architecture patterns are among the technical details and methodologies employed in the papers \cite{b1, b2, b3}. However, the scope of MLLM development varies across the papers, ranging from specific applications to more general frameworks for MLLM development.

While there are no direct contradictions between the papers, some differences in approach and emphasis can be noted. For example, some papers focus on benchmarks and evaluation methodologies \cite{b4, b5}, while others focus on architectural techniques and software engineering practices \cite{b1, b2, b3}. The limitations and challenges acknowledged by the papers include scalability and robustness, cross-modal learning, and the lack of standardization \cite{b1, b4, b5}. Overall, the development of MLLMs is a complex and multidisciplinary field that requires careful consideration of various architectural and technical factors. By highlighting key themes and developments, and discussing implications and connections between different approaches, this survey aims to provide a comprehensive overview of the current state of MLLM development and identify potential avenues for future research.

\section{Vision-Language Pretraining and Masked Language Modeling}
Vision-language pretraining and masked language modeling have emerged as crucial components in the development of multimodal learning models, particularly in the context of Multimodal Learning with Language (MLLM). Recent studies have underscored the significance of these techniques in enhancing model performance across a variety of vision-language tasks. A key finding that prevails across these investigations is the effectiveness of uniform masking in vision-language pretraining \cite{b1}. By increasing the masking rate in Masked Language Modeling (MLM), downstream performance can be significantly improved, and the gap between different masking strategies reduced, thereby making uniform masking competitive with more complex approaches.

Moreover, research has highlighted the importance of efficient data utilization in the cross-modal setting. Traditional MLM can under-utilize images due to short captions or mask stop-words/punctuation, leading to suboptimal fusion of text and image representations \cite{b2}. To address these shortcomings, alternative masking strategies have been proposed that are tailored to the specifics of vision-language pretraining. For instance, Modality Aware Masking (MAM) introduced in MLIM \cite{b3} aims to boost cross-modality interaction by adaptively selecting which modalities to mask during pretraining. Similarly, techniques such as synthesizing hard negatives for Image-Text Matching (ITM) based on language input context have been shown to encourage the model to learn high-quality representations \cite{b4}.

The use of pretraining strategies that enhance cross-modality interaction has also been a focal point of research in this domain. The introduction of additional loss functions, such as image reconstruction (RECON) loss in MLIM \cite{b3}, alongside MLM loss, has been explored to provide a more comprehensive pretraining objective. Furthermore, methods like cross-distillation for MLM have been proposed to improve model robustness by generating soft labels and avoiding the treatment of synonyms as negative samples \cite{b5}. These approaches underscore the importance of developing pretraining strategies that efficiently leverage both visual and textual data.

Despite these advancements, challenges persist in ensuring efficient utilization of cross-modal data and balancing model complexity with performance. The trade-off between simpler models that may benefit from focused pretraining strategies but lack representational capacity, versus more complex architectures, remains a significant consideration \cite{b6}. Additionally, the generalizability of pretraining strategies across various vision-language tasks poses a challenge, as different tasks may require distinct types of cross-modality interactions and representations.

In conclusion, vision-language pretraining and masked language modeling are pivotal for achieving high performance in MLLM. Key themes that emerge from recent research include the importance of efficient data utilization, enhancing cross-modality interaction through tailored pretraining strategies, and addressing the complexities of balancing model simplicity with representational capacity. As the field continues to evolve, future directions may involve exploring more sophisticated masking strategies, integrating additional modalities into pretraining frameworks, and developing more generalized pretraining objectives that can seamlessly adapt across a wide range of vision-language tasks \cite{b7}. The implications of these developments are profound, suggesting potential breakthroughs in areas such as visual question answering, image-text retrieval, and multimodal dialogue systems, thereby underscoring the significance of continued research in this domain.

\section{Applications of MLLMs in Various Domains}
The applications of Multimodal Large Language Models (MLLMs) in various domains have been a subject of increasing interest, with recent studies highlighting their potential to overcome hurdles such as data heterogeneity and sophistication of domain knowledge \cite{b1}. A key finding across these papers is the importance of domain specialization for MLLMs to achieve enhanced performance in specific domains \cite{b2}, \cite{b3}. This is underscored by comprehensive surveys and meta-reviews that systematically categorize existing literature on MLLMs, focusing on their applications, evaluation methodologies, ethical concerns, security, efficiency, and domain-specific applications \cite{b4}, \cite{b5}.

The adaptation of general MLLMs to specific domains through post-training techniques has shown promising results, with advancements in data synthesis, modified training pipelines, and task evaluation contributing to improved performance \cite{b6}. For instance, techniques such as visual instruction synthesizers have been developed to generate diverse visual instruction tasks from domain-specific image-caption pairs, thereby enhancing the domain-specific performance of MLLMs \cite{b7}. Moreover, experiments conducted across different domains, including biomedicine and food, using various MLLM models have demonstrated the effectiveness of domain adaptation techniques \cite{b8}, \cite{b9}.

A recurring theme across these studies is the emphasis on comprehensive evaluation methodologies for MLLMs, recognizing that thorough assessment is crucial for advancing these models \cite{b10}. The community acknowledges the challenges in evaluating MLLMs due to their multifaceted nature, combining various modalities and tasks, and the need for systematic and diverse evaluation methods \cite{b11}. Efforts to develop benchmarks and evaluations for MLLMs are ongoing, with a focus on creating more nuanced and task-specific assessments \cite{b12}.

The integration of multiple modalities in MLLMs is seen as a significant advancement, enabling a more holistic understanding of information that closely mimics human perception \cite{b13}. This development has implications for a wide range of applications, from healthcare and education to entertainment and customer service \cite{b14}. However, despite the progress made, there are identified open challenges and areas for future research, including improving evaluation methods, enhancing domain adaptation techniques, and addressing ethical and security concerns \cite{b15}.

The debate between general versus specialized models is an ongoing one, with some studies advocating for the development of general MLLMs that can be applied across a broad range of tasks, while others argue for the necessity of domain-specialized models to achieve superior performance in specific domains \cite{b16}, \cite{b17}. Furthermore, there are differing approaches to evaluating MLLMs, with some emphasizing the need for more diverse and comprehensive benchmarks, while others focus on developing new evaluation methodologies tailored to the unique capabilities of MLLMs \cite{b18}.

In conclusion, the applications of MLLMs in various domains represent a rapidly evolving field, with key themes including domain specialization, comprehensive evaluation, and multimodality \cite{b19}. As research continues to address the challenges and limitations of these models, including complexity of domain knowledge, evaluation complexity, ethical and security concerns, scalability, and efficiency, we can expect to see significant advancements in their performance and applicability across diverse domains \cite{b20}. Ultimately, the development of effective MLLMs has the potential to transform numerous aspects of our lives, from how we interact with information to how we approach complex problems, highlighting the need for continued investment in this area of research \cite{b21}.

\section{Challenges and Limitations of Current MLLM Systems}
The current state of Multilingual Large Language Models (MLLMs) is marked by significant advancements, but also by several challenges and limitations that hinder their widespread adoption and effective deployment. As identified in \cite{b1}, one of the primary concerns is the lack of systematic approaches to development, evaluation, and deployment, which leads to inefficiencies and redundancies throughout the MLLM ecosystem. This issue is further complicated by the rapid evolution of MLLMs, making it challenging to establish stable benchmarks and best practices \cite{b2}. The absence of standardized frameworks for MLLM development, deployment, and evaluation exacerbates this problem, as noted in \cite{b3}, highlighting the need for a principled process like the Machine Learning Technology Readiness Levels (MLTRL) framework \cite{b4} to ensure robust, reliable, and responsible ML systems.

Another critical challenge facing MLLMs is ensuring reliability and robustness, particularly in critical applications where certification is required \cite{b5}. The integration of ML techniques into certified systems poses significant risks due to the inherent uncertainties associated with these techniques. Addressing this challenge requires not only technical advancements but also new engineering practices that can mitigate these risks \cite{b6}. Furthermore, the deployment and scaling of MLLMs are fraught with difficulties, as discussed in \cite{b7}, which emphasizes the need for innovations in serving systems to optimize performance and efficiency without compromising on core decoding mechanisms.

The issue of redundancy, especially among MLLM benchmarks, is a significant concern that affects the effectiveness of evaluations \cite{b8}. Proposals such as redundancy principles for MLLM benchmarks aim to refine these evaluations by eliminating unnecessary redundancies and promoting more comprehensive analyses of existing benchmarks \cite{b9}. This effort towards efficiency and effectiveness is crucial for advancing the field, as it enables a more accurate assessment of MLLM capabilities and identifies areas that require further research and development.

In addition to these technical challenges, there are also broader systemic issues at play. The contrast between focusing on technical versus systemic challenges \cite{b10} and the approach of proposing specific solutions versus providing comprehensive analyses \cite{b11} reflects the complexity and multifaceted nature of the problems facing MLLMs. Some studies focus on general ML challenges relevant to a wide range of applications, while others delve into issues unique to MLLMs, such as benchmark redundancy and the specifics of LLM serving systems \cite{b12}. This diversity in approaches and focuses underscores the need for a holistic understanding of MLLM challenges and limitations.

The implications of these challenges are far-reaching, affecting not only the development and deployment of MLLMs but also their potential applications and societal impact. For instance, the reliability and robustness of MLLMs are critical in applications such as language translation, sentiment analysis, and text generation, where inaccuracies or biases can have significant consequences \cite{b13}. Moreover, the lack of standardization and the presence of redundancy in benchmarks can hinder the progress of research in this area, leading to inefficiencies in resource allocation and a slower pace of innovation \cite{b14}.

In conclusion, the challenges and limitations of current MLLM systems are multifaceted and interconnected, requiring a comprehensive and systematic approach to address them effectively. By understanding these challenges and their implications, researchers and practitioners can work towards developing more effective, efficient, and reliable MLLM systems that can fully realize their potential in various applications and contribute positively to society. Future research should focus on developing standardized frameworks for MLLM development and evaluation, enhancing reliability and robustness, especially in certified systems, and addressing the inefficiencies and redundancies that currently plague the field \cite{b15}.

\section{Hallucination and Toxicity Detection in MLLMs}
Hallucination and toxicity detection in Multimodal Large Language Models (MLLMs) has emerged as a critical area of research, with significant implications for the reliability and trustworthiness of these models in various applications, including mental health support and medical contexts. Recent studies have proposed innovative approaches to address this issue, highlighting the need for effective detection and mitigation strategies \cite{b1}, \cite{b2}. A comprehensive review of the current literature on hallucination detection and mitigation underscores the importance of developing more robust methods to tackle this problem \cite{b3}. 

The introduction of novel meta-evaluation benchmarks, such as MHaluBench, and unified multimodal hallucination detection frameworks, like UNIHD, has facilitated the evaluation of advancements in hallucination detection methods \cite{b4}. Furthermore, the creation of domain-specific benchmarks, including Med-HallMark, the first benchmark designed for hallucination detection and evaluation within the medical multimodal domain, has emphasized the need to address domain-specific requirements and constraints \cite{b5}. The proposal of new evaluative metrics, such as the MediHall Score, has also contributed to a more comprehensive assessment of MLLMs' hallucinations in medical contexts.

A common theme across these studies is the emphasis on multimodal approaches, which leverage ensemble methods, knowledge graphs, and vector stores to enhance hallucination detection and mitigation \cite{b2}, \cite{b4}. The use of self-reflection schemes and multi-agent debate approaches has also been proposed to promote slow-thinking and divergent-thinking, thereby mitigating hallucinations in MLLMs \cite{b1}. However, the papers also highlight the challenges associated with scalability and generalizability, as well as the need for robust evaluation metrics and benchmarks \cite{b3}, \cite{b4}.

While some studies focus primarily on detecting hallucinations \cite{b3}, \cite{b4}, others emphasize the importance of mitigating hallucinations \cite{b1}, \cite{b2}. The choice between single-task and multi-task approaches is also a point of discussion, with some papers proposing single-task methods \cite{b3} and others suggesting multi-task approaches that leverage auxiliary tools and ensemble methods \cite{b4}, \cite{b5}. 

The technical details and methodologies employed in these studies are diverse, ranging from self-reflection schemes and ensemble methods to meta-evaluation benchmarks and multitask training approaches \cite{b1}, \cite{b2}, \cite{b4}, \cite{b5}. Despite the progress made, the papers acknowledge several limitations and challenges, including the lack of effective detection and mitigation strategies, scalability and generalizability issues, and the need for more robust evaluation metrics and benchmarks \cite{b1}, \cite{b2}, \cite{b3}, \cite{b4}, \cite{b5}. 

Overall, the research on hallucination and toxicity detection in MLLMs highlights the complexity of this issue and the need for continued innovation and development of more effective methods to address it. As MLLMs become increasingly ubiquitous in various applications, the importance of ensuring their reliability and trustworthiness cannot be overstated, underscoring the significance of this research area and its potential implications for the future of natural language processing and multimodal interaction.

\section{Model Compression, Acceleration, and Efficient Training of LLMs}
Model compression, acceleration, and efficient training of Large Language Models (LLMs) have become crucial aspects of research in the field of natural language processing \cite{b1}. The increasing complexity and size of LLMs have led to significant computational requirements, making it essential to develop techniques that can reduce their computational footprint while maintaining performance \cite{b2}. Recent studies have introduced various model compression techniques, including quantization, pruning, and knowledge distillation, which have shown promising results in reducing the computational requirements of LLMs \cite{b3}. For instance, the paper "Continuous Approximations for Improving Quantization Aware Training of LLMs" presents two continuous approximations to the quantization aware training (QAT) process, which improves the performance of quantized models and reduces energy consumption \cite{b4}.

The importance of model compression is further emphasized by surveys such as "A Survey on Model Compression for Large Language Models" and "Faster and Lighter LLMs: A Survey on Current Challenges and Way Forward", which provide a comprehensive overview of model compression techniques \cite{b5}, \cite{b6}. These surveys highlight the need for efficient training methods, including optimization strategies and learning acceleration techniques, to reduce the time and resources required for training LLMs. Optimization strategies, such as skipping latter attention sublayers in Transformer LLMs, have been explored in papers like "Enhancing Inference Efficiency of Large Language Models: Investigating Optimization Strategies and Architectural Innovations" \cite{b7}. Additionally, learning acceleration techniques, such as the theory for optimal learning of language models presented in "Towards Optimal Learning of Language Models", aim to reduce the necessary training steps for achieving superior performance \cite{b8}.

A common theme among these studies is the trade-off between performance and efficiency. While model compression techniques can significantly reduce computational requirements, they often come at the cost of performance degradation \cite{b9}. Therefore, it is essential to balance model accuracy with computational requirements and energy consumption. The papers analyzed in this section employ a range of technical details and methodologies, including quantization aware training, Transformer architectures, and benchmarking datasets like perplexity (PPL) and accuracy \cite{b10}. However, the field still faces significant challenges, such as scalability and energy efficiency, which must be addressed to ensure sustainable AI development \cite{b11}.

The differences in focus and methodology among the analyzed papers are noteworthy. While some papers focus on quantization as a primary compression technique, others emphasize pruning or knowledge distillation \cite{b12}. Furthermore, optimization strategies can be explored at both the architectural and model levels, as seen in papers like "Enhancing Inference Efficiency of Large Language Models" \cite{b7}. Despite these differences, the overall goal of developing efficient and scalable LLMs remains a common objective. The implications of these developments are significant, as they have the potential to enable the widespread adoption of LLMs in various applications, from natural language processing to computer vision \cite{b13}. Ultimately, the connections between model compression, acceleration, and efficient training highlight the need for a holistic approach to developing sustainable and efficient AI systems.

\section{Evaluating MLLMs: Benchmarks, Metrics, and Adversarial Testing}
Evaluating Multimodal Large Language Models (MLLMs) is a crucial aspect of their development, as it enables researchers to assess their performance, identify areas for improvement, and ensure their reliability and trustworthiness \cite{b1}. The rapid advancement of MLLMs has led to an increased need for comprehensive evaluation methods, prompting researchers to investigate various benchmarks, metrics, and adversarial testing techniques \cite{b2}. This section provides an overview of the current state of evaluating MLLMs, highlighting key findings, common themes, and technical details.

The issue of redundancy in MLLM benchmarks has been highlighted by several studies, including "Redundancy Principles for MLLMs Benchmarks" \cite{b3}, which proposes principles to address this problem and improve benchmark construction. Comprehensive surveys of MLLM evaluation, such as "A Survey on Benchmarks of Multimodal Large Language Models" \cite{b4}, "Surveying the MLLM Landscape: A Meta-Review of Current Surveys" \cite{b5}, and "MME-Survey: A Comprehensive Survey on Evaluation of Multimodal LLMs" \cite{b6}, provide an overview of existing benchmarks, evaluation methods, and challenges in MLLM research. These studies emphasize the need for robust evaluation methods that consider the complexity and applications of MLLMs.

Evaluation techniques for trustworthiness, such as those introduced in "Enhancing Trust in LLMs: Algorithms for Comparing and Interpreting LLMs" \cite{b7}, are crucial for assessing MLLM performance. These techniques include perplexity measurement, NLP metrics, and adversarial testing, which help ensure the reliability and fairness of MLLMs. Moreover, several studies stress the importance of human evaluation in capturing nuances that automated metrics may miss \cite{b8}. Human evaluation requires significant expertise and resources, which can be a limitation in large-scale MLLM development.

A common theme among these studies is the need for comprehensive evaluation methods that consider the increasing complexity and applications of MLLMs. Redundancy and inefficiency in benchmark construction are also highlighted as significant issues, with proposals for strategies to address these problems \cite{b3}. Transparency and trustworthiness are emphasized as essential aspects of MLLM evaluation, ensuring reliable and fair development \cite{b7}.

Contrasting viewpoints on quantitative vs. qualitative evaluation methods are evident in the literature. While some studies focus on quantitative metrics \cite{b4}, others emphasize the importance of human evaluation and qualitative assessment \cite{b8}. Different approaches to benchmark construction are also proposed, ranging from redundancy reduction to comprehensive surveying of existing benchmarks \cite{b3} \cite{b5}.

Technical details and methodologies for MLLM benchmark construction, evaluation metrics, and algorithms are discussed in various studies. For example, "Enhancing Trust in LLMs: Algorithms for Comparing and Interpreting LLMs" \cite{b7} introduces a range of evaluation metrics and algorithms, including perplexity measurement, NLP metrics, and adversarial testing. Stratified evaluation methods are also proposed to assess MLLM performance across different capabilities and applications \cite{b6}.

Despite the progress made in evaluating MLLMs, several limitations and challenges remain. Redundancy and inefficiency in benchmark construction can lead to redundant and ineffective evaluation methods \cite{b3}. The increasing complexity and scale of MLLMs pose significant challenges for developing comprehensive evaluation methods \cite{b2}. Moreover, human evaluation requires significant expertise and resources, which can be a limitation in large-scale MLLM development \cite{b8}.

In conclusion, evaluating MLLMs is a crucial aspect of their development, requiring comprehensive and robust evaluation methods. By analyzing the current state of MLLM evaluation, we can identify key areas for future research and development, including the creation of more efficient and effective benchmark construction strategies, the development of comprehensive evaluation methods, and the integration of human evaluation and expertise \cite{b1} \cite{b7}. Ultimately, addressing these challenges will enable researchers to develop reliable, trustworthy, and high-performing MLLMs that can be applied in various domains.

\section{Cognitive Abilities and Understanding in Multimodal Intelligence}
The cognitive abilities and understanding in Multimodal Large Language Models (MLLMs) have been a subject of extensive research, with recent studies shedding light on their capabilities and limitations \cite{b1}. A significant contribution to this field is the development of comprehensive evaluation benchmarks, such as the Multimodal Large Language Model Evaluation benchmark (MME) and CogDevelop2K, which assess both perception and cognition abilities across various subtasks \cite{b2}. These benchmarks have facilitated a deeper understanding of MLLMs' strengths and weaknesses, highlighting their impressive performance on complex tasks while also revealing a reversed cognitive developmental trajectory compared to humans \cite{b3}. This phenomenon, observed in CogDevelop2K, suggests that while MLLMs can excel in certain areas, their foundational understanding and cognitive development may not mirror human intelligence \cite{b4}.

The importance of evaluation as a discipline for advancing the field of MLLMs cannot be overstated, with a strong emphasis on systematic reviews of existing evaluation methods and the proposal of new benchmarks to assess multimodal reasoning, perception, cognition, and trustworthiness \cite{b5}. A common theme emerging from recent studies is the exploration of MLLMs' capabilities in multimodal reasoning and intelligence, which is seen as a pathway towards achieving Artificial General Intelligence (AGI) \cite{b6}. However, there is a recurring pattern highlighting the necessity for comprehensive evaluation frameworks that can accurately assess the wide range of capabilities of MLLMs, including perception, cognition, and domain-specific applications \cite{b7}.

The approach to assessing cognitive abilities in MLLMs varies, with some studies emphasizing the importance of understanding foundational cognitive concepts, such as object permanence, and others focusing on more complex reasoning tasks \cite{b8}. Technical details, such as the design of new benchmarks like MME and CogDevelop2K, provide valuable insights into the complexity of assessing multimodal intelligence \cite{b9}. The use of various metrics for evaluating MLLMs, including those that measure perception, cognition, reasoning, and trustworthiness, underscores the multifaceted nature of intelligence \cite{b10}.

Despite the impressive performances of MLLMs, significant room for improvement remains in terms of genuine understanding and cognitive abilities \cite{b11}. The current evaluation methods for MLLMs are recognized as having limitations, suggesting a need for further research into developing more effective assessment frameworks \cite{b12}. Interpreting the results of evaluations, especially in relation to human-like intelligence and cognitive development, poses a challenge due to the complex and multi-faceted nature of intelligence \cite{b13}. Addressing these challenges will be crucial for advancing the field towards more sophisticated models of artificial intelligence.

The analysis of recent studies highlights key themes and developments in the field of MLLMs, including the importance of comprehensive evaluation frameworks, the potential of MLLMs for multimodal reasoning and intelligence, and the need for further research into developing more effective assessment frameworks \cite{b14}. The implications of these findings are significant, with connections to broader discussions around the development of artificial intelligence and its potential applications \cite{b15}. As research in this field continues to evolve, it is likely that we will see significant advancements in our understanding of MLLMs' cognitive abilities and their potential for achieving true artificial intelligence.

\section{Future Directions and Open Challenges in MLLM Research}
The future directions and open challenges in Multimodal Large Language Model (MLLM) research are multifaceted and far-reaching, reflecting the rapid evolution of this field and its potential to transform various domains. As highlighted by recent surveys and systematic mapping studies \cite{b1}, \cite{b2}, a key area that requires further investigation is the development of more effective evaluation methods for MLLMs. The need for better evaluation benchmarks and methodologies is widely acknowledged, as current approaches often focus on performance metrics that do not necessarily capture the complexity and nuances of real-world applications \cite{b3}. This limitation underscores the importance of connecting MLLM research to societal and scientific problems, rather than solely pursuing incremental improvements in performance metrics \cite{b4}.

The connection to real-world problems is a recurring theme in MLLM research, with papers emphasizing the need for machine learning research to focus on problems that matter to society and science \cite{b5}. This includes applications in hypothesis discovery, experiment planning, scientific writing, and peer reviewing, where MLLMs have shown significant potential \cite{b6}. Furthermore, the growing interest in multimodal MLLMs is evident, with applications in visual question answering, perception, understanding, and reasoning \cite{b7}. However, developing effective multimodal MLLMs poses significant challenges, including the integration of multiple modalities and the need for more sophisticated evaluation methods \cite{b8}.

Another critical aspect of MLLM research is the importance of testing and validation techniques. Researchers stress the need for rigorous testing and validation to ensure the reliability and trustworthiness of MLLMs \cite{b9}. This includes the development of non-functional testing and reinforcement learning-based testing techniques, which have shown promise in evaluating the performance of MLLMs \cite{b10}. Nevertheless, the lack of empirical evidence to support the effectiveness of various evaluation methods and testing techniques is a significant limitation that needs to be addressed through further research \cite{b11}.

The technical details and methodologies employed in MLLM research are also worthy of note. Systematic mapping studies have provided comprehensive overviews of the current state of MLLM research, identifying trends, themes, and areas for further investigation \cite{b12}. Additionally, researchers have proposed various benchmarks and evaluation metrics to assess the performance of MLLMs, including perception and understanding, cognition and reasoning, and specific domain-based evaluations \cite{b13}. However, the limited availability of publicly available tools and resources remains a significant challenge, hindering the development and testing of MLLMs \cite{b14}.

In conclusion, the future directions and open challenges in MLLM research are complex and multifaceted. Addressing these challenges will require a concerted effort from researchers to develop more effective evaluation methods, connect MLLM research to real-world problems, and improve testing and validation techniques. Furthermore, the development of multimodal MLLMs and the integration of multiple modalities pose significant technical challenges that need to be overcome. By acknowledging these limitations and challenges, researchers can work towards creating more robust, reliable, and impactful MLLMs that have the potential to transform various domains and drive scientific progress \cite{b15}. Ultimately, the implications of MLLM research are far-reaching, with potential connections to various fields, including machine learning, scientific research, and multimodal applications, highlighting the need for continued innovation and advancement in this field \cite{b16}.

\section{Ethical Considerations and Implications of MLLMs for Society}
The development and deployment of Machine Learning Large Language Models (MLLMs) have significant ethical considerations and implications for society, emphasizing the need for a thorough examination of their potential consequences \cite{b1}. A comprehensive analysis of relevant papers reveals that integrating ethical standards into the development and deployment of MLLMs is crucial to ensure responsible innovation \cite{b2}. The importance of evaluating and addressing potential biases, fairness issues, and other ethical concerns in MLLMs cannot be overstated, as these models can have a profound impact on society \cite{b3}. 

Researchers have proposed various frameworks and guidelines for ensuring high ethical standards in MLLM development, including conceptual and practical principles for ML research and deployment \cite{b4}, comprehensive surveys of ethical challenges associated with Large Language Models (LLMs) \cite{b5}, and utility frameworks for characterizing the trade-off between information gain and potential ethical harms in ML evaluation \cite{b6}. These efforts highlight the complexity of integrating ethical considerations into MLLM development and deployment, as well as the need for ongoing research into the ethical implications of these models \cite{b7}.

A key theme emerging from the analysis is the recognition that MLLMs can have significant societal impact and must be designed and deployed responsibly \cite{b8}. This requires not only a deep understanding of the technical aspects of MLLMs but also a nuanced appreciation of their potential consequences on individuals, communities, and society as a whole \cite{b9}. Furthermore, the development of practical guidelines and frameworks for ensuring ethical standards in MLLM research and industry applications is essential to mitigate potential risks and harms \cite{b10}.

The analysis also reveals contrasting viewpoints and approaches to addressing the ethical implications of MLLMs. While some researchers focus on integrating ethical considerations into the development of MLLMs, others emphasize the need for evaluating and addressing potential biases and fairness issues in existing systems \cite{b11}. Additionally, there are differing perspectives on the ability of MLLMs to navigate ethical dilemmas and make tough choices, with some studies suggesting that these models may struggle with complex moral decisions \cite{b12}.

Technical details and methodologies employed in the analysis include surveys, evaluations, and conceptual frameworks \cite{b13}. For example, one study used a dataset of 1,730 ethical dilemmas to evaluate the performance of 20 well-known LLMs \cite{b14}. These technical approaches highlight the complexity of evaluating and addressing ethical concerns in MLLMs and the need for ongoing research into the development of more effective methodologies.

Despite the progress made in understanding the ethical implications of MLLMs, several limitations and challenges remain. These include the complexity of integrating ethical considerations into MLLM development and deployment, the need for ongoing research into the ethical implications of these models, and the challenge of evaluating and addressing potential biases and fairness issues in existing systems \cite{b15}. Moreover, the development of practical guidelines and frameworks that can be applied across different contexts and applications is a significant challenge \cite{b16}. Ultimately, addressing these limitations and challenges will require a collaborative effort from researchers, policymakers, and industry stakeholders to ensure that MLLMs are developed and deployed in a responsible and ethical manner \cite{b17}.

\section{Technical Advances and Emerging Trends in MLLM Development}
The development of Multimodal Large Language Models (MLLMs) has witnessed significant technical advances and emerging trends in recent years. A comprehensive analysis of relevant papers reveals several key findings and contributions that have shaped the field. Recent advancements in LLM serving systems, for instance, focus on system-level enhancements that improve performance and efficiency without altering core decoding mechanisms \cite{b1}. This is complemented by a thorough understanding of Machine Learning Operations (MLOps) that covers the entire machine learning life cycle, providing a clear picture of the state-of-the-art ML technologies \cite{b2}. The importance of multimodal benchmarks in evaluating MLLMs is also highlighted, with a review of 211 benchmarks across four core domains: understanding, reasoning, generation, and application \cite{b3}, and another review of 200 benchmarks emphasizing the role of perception, cognition, and reasoning in multimodal understanding \cite{b4}. Furthermore, the introduction of Machine Learning Technology Readiness Levels (MLTRL) offers a principled process for developing and deploying robust, reliable, and responsible ML systems \cite{b5}.

A common theme that emerges across these papers is the emphasis on evaluation and benchmarking, with multiple studies underscoring the need for standardized methods to assess MLLM performance \cite{b3, b4}. System-level enhancements are another recurring theme, focusing on efficiency, scalability, and reliability \cite{b1}. The consideration of the entire machine learning life cycle, from research to deployment, is also a dominant theme in MLLM development, highlighting the importance of MLOps and ML development \cite{b2, b5}. Additionally, the ability of MLLMs to understand and generate multimodal content is a significant area of exploration, with various papers delving into perception, cognition, and reasoning in this context \cite{b3, b4}.

While there is a consensus on these themes, some contrasting viewpoints or approaches are notable. For example, some papers focus on system-level enhancements \cite{b1, b5}, whereas others emphasize the importance of model architectures and benchmarking \cite{b3, b4}. There is also a distinction between MLOps as a distinct discipline \cite{b2} and its parallels with traditional software engineering \cite{b5}. From a technical standpoint, the papers vary in their methodologies, including discussions on LLM serving systems \cite{b1}, detailed analyses of multimodal benchmarks \cite{b3, b4}, and the introduction of the MLTRL framework \cite{b5}.

Despite these advances, several limitations and challenges persist in MLLM development. A significant concern is the lack of standardized evaluation methods for MLLMs, which current studies note as a limitation \cite{b3, b4}. There are also warnings against rushing ML development, given the risks of technical debt, scope creep, and model misuse \cite{b5}. The need for robust and reliable systems that can efficiently serve multimodal content is another critical challenge, requiring careful consideration of system-level enhancements, evaluation, and deployment. These challenges underscore the complexity and multidisciplinary nature of MLLM development, necessitating continued research and innovation to address the technical, methodological, and ethical implications of these models. Ultimately, the future of MLLMs depends on overcoming these challenges through collaborative efforts that integrate advances in model architectures, system-level enhancements, and rigorous evaluation methodologies.

\section{Case Studies: Real-World Applications and Success Stories of MLLMs}
The case studies and real-world applications of Multitask Large Language Models (MLLMs) demonstrate their potential to transform various industries and domains. A systematic analysis of relevant papers reveals key findings and contributions, including the challenges and limitations associated with MLLMs \cite{b1}. For instance, the study "No Size Fits All: The Perils and Pitfalls of Leveraging LLMs Vary with Company Size" highlights the importance of considering company size when deploying MLLMs, providing a practical guide for industries to utilize these models more efficiently \cite{b2}. Another paper, "A Primer on Large Language Models and their Limitations," offers an introduction to MLLMs, their strengths, limitations, applications, and research directions, making it a valuable resource for those new to the field \cite{b3}.

The analysis of these papers also reveals common themes and patterns, such as the need for practical guides and resources to facilitate the adoption and utilization of MLLMs \cite{b4}. The paper "Towards a Middleware for Large Language Models" proposes a middleware system architecture to address challenges like complexity and integration issues with existing systems \cite{b5}. Furthermore, the study "A Systematic Mapping Study on Testing of Machine Learning Programs" provides an overview of the area of testing ML programs, identifying trends and highlighting the need for more empirical evidence and publicly available tools \cite{b6}.

The contrasting viewpoints and approaches presented in these papers demonstrate the diversity of research in MLLMs. For example, some papers focus on industrial applications and challenges \cite{b2}, while others provide a more academic introduction to the field \cite{b3}. Technical approaches, such as the proposal of a middleware system architecture \cite{b5}, are also presented alongside more methodological approaches, like systematic mapping studies \cite{b6}. Additionally, some papers emphasize the challenges and limitations of MLLMs \cite{b1}, while others focus on success stories and applications \cite{b2}.

The technical details and methodologies employed in these papers are also noteworthy. Systematic mapping studies \cite{b6} and case studies \cite{b2} are used to identify and analyze existing literature, providing valuable insights into the challenges and successes of MLLM deployment. The proposal of a middleware system architecture \cite{b5} demonstrates the technical solutions being developed to address the complexity and integration issues associated with MLLMs.

Despite the advancements in MLLMs, several limitations and challenges remain. Complexity and integration issues are significant concerns, as these models require substantial resources and expertise to deploy and integrate with existing systems \cite{b1}. The need for more empirical evidence to compare and assess the effectiveness of different techniques and approaches in testing ML programs is also a pressing issue \cite{b6}. Furthermore, scalability and customization can be difficult, particularly for smaller organizations or those with limited resources \cite{b2}. Finally, the use of MLLMs raises concerns about data privacy and security, particularly when deployed in cloud-based environments \cite{b5}.

In conclusion, the case studies and real-world applications of MLLMs demonstrate their potential to transform various industries and domains. However, it is essential to address the challenges and limitations associated with these models, including complexity, integration issues, scalability, customization, and privacy and security concerns. By acknowledging these limitations and continuing to develop practical guides, resources, and technical solutions, researchers and practitioners can unlock the full potential of MLLMs and drive innovation in this rapidly evolving field \cite{b1}, \cite{b2}, \cite{b3}, \cite{b4}, \cite{b5}, \cite{b6}.

\section{Conclusion and Outlook for the Future of Multimodal AI Systems}
The conclusion and outlook for the future of multimodal AI systems, particularly Multimodal Large Language Models (MLLMs), are promising, with significant contributions made to the field in recent years \cite{b1}. The potential of MLLMs to achieve artificial general intelligence (AGI) has been highlighted through their impressive capabilities across a wide range of multimodal tasks and applications, including natural language, vision, audio, and physiological sequences \cite{b2}, \cite{b3}. Evaluation methods for MLLMs have also been developed to assess their capabilities across different dimensions, such as perception, reasoning, trustworthiness, and domain-specific applications \cite{b4}, emphasizing the importance of regarding evaluation as a critical discipline essential for advancing the field of MLLMs.

A common theme among researchers is the emphasis on integrating multiple modalities to enable more comprehensive understanding and generation of multimodal content \cite{b5}. This integration is crucial for achieving AGI, as it allows MLLMs to learn from diverse sources of data and improve their performance across various tasks. The development of evaluation methods and benchmarks has also been a key focus area, providing a comprehensive overview of the current state of MLLM evaluation \cite{b6}. Furthermore, the importance of reasoning abilities in MLLMs, including abstract reasoning and multimodal reasoning, has been recognized as a critical aspect of achieving AGI \cite{b7}.

Despite the progress made, there are contrasting viewpoints and approaches among researchers. For instance, some studies focus on specific modalities, such as vision or audio, while others adopt a more integrated approach that combines multiple modalities \cite{b8}. Additionally, variations in evaluation methods and benchmarks have been observed, with some studies emphasizing perception and understanding, while others focus on cognition and reasoning \cite{b9}. These differences highlight the need for more diverse and robust evaluation datasets and metrics to support the development of more capable and reliable MLLMs.

The technical methodologies employed in MLLM research are diverse, ranging from comprehensive reviews of existing evaluation protocols and benchmarks to systematic analyses of task designs, evaluation metrics, and dataset constructions across diverse modalities \cite{b10}. Investigations into the reasoning abilities of MLLMs have also been conducted, including abstract reasoning and multimodal reasoning \cite{b11}. However, limitations and challenges persist, such as the need for more robust and diverse evaluation datasets and metrics, addressing the shortcomings of current MLLMs, and developing evaluation methods that can effectively assess the complex and multifaceted capabilities of MLLMs \cite{b12}.

In conclusion, the future of multimodal AI systems, particularly MLLMs, holds significant promise, with potential applications in various domains. To realize this potential, it is essential to address the limitations and challenges identified in this survey, including the development of more robust and diverse evaluation datasets and metrics, improving the reasoning abilities of MLLMs, and investigating the implications and connections of MLLMs in real-world settings \cite{b13}. As researchers continue to advance the field of MLLMs, it is crucial to maintain a focus on achieving AGI, while also ensuring that these models are reliable, trustworthy, and aligned with human values \cite{b14}. Ultimately, the successful development of MLLMs has the potential to transform numerous aspects of society, from healthcare and education to entertainment and commerce, making it an exciting and worthwhile area of research \cite{b15}.

\section{Conclusion}
The analysis of papers related to Multimodal Large Language Models (MLLMs) reveals several key findings, common themes, and limitations. The development of comprehensive evaluation benchmarks for MLLMs is crucial for advancing the field, as highlighted in papers such as \cite{b1} and \cite{b2}. These benchmarks are essential for assessing the performance of MLLMs in various applications, including visual question answering, visual perception, understanding, and reasoning, as demonstrated in papers like \cite{b3} and \cite{b4}. The importance of standardized evaluation frameworks for MLLMs is emphasized in papers such as \cite{b5}, which proposes a holistic framework for evaluating MLLMs. This emphasis on standardization is a common theme across various papers, including \cite{b1} and \cite{b2}, highlighting the need for a unified approach to evaluating MLLMs.

The rapid evolution of MLLMs has led to a growing need for comprehensive evaluation benchmarks, as noted in multiple papers, including \cite{b3} and \cite{b4}. The importance of multimodal understanding and reasoning is another common theme, with papers like \cite{b1} and \cite{b2} proposing different approaches to evaluating these aspects. However, the development of evaluation benchmarks and frameworks involves careful consideration of task design, dataset construction, and metric selection, as discussed in papers like \cite{b5} and \cite{b6}. Various methodologies are employed to evaluate MLLMs, including quantitative statistics and qualitative analysis, as seen in papers such as \cite{b1} and \cite{b7}.

Despite the progress made in MLLM research, several limitations and challenges remain. The lack of standardized evaluation frameworks hinders the comparison and evaluation of different MLLM models, as highlighted in papers such as \cite{b5}. Further research is needed to address the limitations and challenges of MLLMs, including issues related to calibration, in-context learning, instruction following, language performance, hallucination, and robustness, as discussed in papers like \cite{b1} and \cite{b8}. The recognition of these limitations and challenges is a common thread throughout the papers, with many authors highlighting the need for further investigation. Different papers propose varying evaluation benchmarks and frameworks, reflecting diverse approaches to evaluating MLLMs. For example, \cite{b2} proposes a comprehensive evaluation benchmark, while \cite{b5} presents a standardized framework for assessing MLLMs.

The implications of these findings are significant, as they highlight the need for continued research and development in MLLM evaluation. The connections between different papers and themes are also noteworthy, as they demonstrate the evolving nature of MLLM research. For instance, the emphasis on multimodal understanding and reasoning in \cite{b1} is echoed in \cite{b2}, which proposes a framework for evaluating these aspects. Similarly, the discussion of limitations and challenges in \cite{b5} is complemented by the proposals for addressing these issues in \cite{b8}. Ultimately, the analysis of papers related to MLLMs reveals a complex and dynamic field, with many opportunities for future research and development. As the field continues to evolve, it is essential to prioritize the development of comprehensive evaluation benchmarks and standardized frameworks, as well as to address the limitations and challenges associated with MLLMs, in order to fully realize their potential.

\end{document}